\begin{lemma}
\label{lemma-fs-local}
Let $(R, \mathfrak m) \to (S, \mathfrak n)$ be a local homomorphism
of local rings. The following are equivalent
\begin{enumerate}
\item $R \to S$ is formally smooth in the $\mathfrak n$-adic topology,
\item for every solid commutative diagram
$$
\xymatrix{
S \ar[r] \ar@{-->}[rd] & A/J \\
R \ar[r] \ar[u] & A \ar[u]
}
$$
of local homomorphisms of local rings where $J \subset A$ is
an ideal of square zero, $\mathfrak m_A^n = 0$ for some $n > 0$, and
$S \to A/J$ induces an isomorphism on residue fields, a dotted
arrow exists which makes the diagram commute.
\end{enumerate}
If $S$ is Noetherian these conditions are also equivalent to
\begin{enumerate}
\item[(3)] same as in (2) but only for diagrams where in addition
$A \to A/J$ is a small extension
(Algebra, Definition \ref{algebra-definition-small-extension}).
\end{enumerate}
\end{lemma}

\begin{proof}
The implication (1) $\Rightarrow$ (2) follows from the definitions.
Consider a diagram
$$
\xymatrix{
S \ar[r] \ar@{-->}[rd] & A/J \\
R \ar[r] \ar[u] & A \ar[u]
}
$$
as in Definition \ref{definition-formally-smooth} for the
$\mathfrak m$-adic topology on $R$ and the $\mathfrak n$-adic topology on $S$.
Pick $m > 0$ with $\mathfrak n^m(A/J) = 0$
(possible by continuity of maps in diagram).
Consider the subring $A'$ of $A$ which is the inverse image
of the image of $S$ in $A/J$. Set $J' = J$ viewed as an ideal in $A'$.
Then $J'$ is an ideal of square zero in $A'$ and $A'/J'$
is a quotient of $S/\mathfrak n^m$. Hence $A'$ is local
and $\mathfrak m_{A'}^{2m} = 0$. Thus we get a diagram
$$
\xymatrix{
S \ar[r] \ar@{-->}[rd] & A'/J' \\
R \ar[r] \ar[u] & A' \ar[u]
}
$$
as in (2). If we can construct the dotted arrow in this diagram,
then we obtain the dotted arrow in the original one by composing
with $A' \to A$. In this way we see that (2) implies (1).

\medskip\noindent
Assume $S$ Noetherian. The implication (1) $\Rightarrow$ (3) is immediate.
Assume (3) and suppose a diagram as in (2) is given.
Then $\mathfrak m_A^n J = 0$ for some $n > 0$.
Considering the maps
$$
A \to A/\mathfrak m_A^{n - 1}J \to \ldots \to A/\mathfrak mJ \to A/J
$$
we see that it suffices to produce the lifting if $\mathfrak m_A J = 0$.
Assume $\mathfrak m_A J = 0$ and let $A' \subset A$ be the ring
constructed above. Then $A'/J'$ is Artinian
as a quotient of the Artinian local ring $S/\mathfrak n^m$.
Thus it suffices to show that given property (3) we can find the dotted
arrow in diagrams as in (2) with $A/J$ Artinian and
$\mathfrak m_A J = 0$.
Let $\kappa$ be the common residue field of $A$,
$A/J$, and $S$. By (3), if $J_0 \subset J$ is an ideal with
$\dim_\kappa(J/J_0) = 1$, then we can produce a dotted arrow
$S \to A/J_0$. Taking the product we obtain
$$
S \longrightarrow \prod\nolimits_{J_0 \text{ as above}} A/J_0
$$
Clearly the image of this arrow is contained in the sub $R$-algebra
$A'$ of elements which map into the small diagonal
$A/J \subset \prod_{J_0} A/J$.
Let $J' \subset A'$ be the elements mapping to zero in $A/J$.
Then $J'$ is an ideal of square zero and as $\kappa$-vector space
equal to
$$
J' = \prod\nolimits_{J_0 \text{ as above}} J/J_0
$$
Thus the map $J \to J'$ is injective. By the theory of vector spaces
we can choose a splitting $J' = J \oplus M$. It follows that
$$
A' = A \oplus M
$$
as an $R$-algebra. Hence the map $S \to A'$ can be composed with
the projection $A' \to A$ to give the desired dotted arrow thereby
finishing the proof of the lemma.
\end{proof}